% !TeX root = ../../Book1.tex
\chapter{Sobolev 空间中的紧性}
\label{b1:ch:25}
\BookChapterNumber{b1:ch:25}
\begin{chapterabstract}
本章把 Sobolev 范数的统一估计转化为函数列的收敛结论。先以有限网格平均证明有限指数 \(L^p\) 的平移紧性判据，再结合延拓、梯度平移估计及插值得到 Rellich–Kondrachov 定理。随后区分弱极限的识别、较弱范数中的强收敛和原 Sobolev 范数的强收敛，并在具体变分问题中说明各自的用途。临界集中和粗糙边界的例子用来检验结论的适用范围。
\end{chapterabstract}

\begin{learninggoals}
\begin{enumerate}
\item\label{b1:goal:25-fk}用尾部和一致平移连续性构造有限维逼近，证明 Fréchet–Kolmogorov 判据，辨别单个函数的连续性与函数族的一致条件。
\item\label{b1:goal:25-rellich}从弱梯度控制小平移，证明积分范数和连续函数范数中的紧嵌入，并说明严格低于临界指数的必要性。
\item\label{b1:goal:25-boundary}区分一般 Lipschitz 区域、任意有界开集上的零边界闭包及内部紧性，高阶结论保留相应的延拓假设。
\item\label{b1:goal:25-weak}识别 Sobolev 弱极限，分别通过紧映射和标准范数的收敛恢复强收敛，避免对非线性表达式直接代入弱极限。
\item\label{b1:goal:25-variation}完整实施一个凸能量的极小值证明，并在习题中用紧性保留非凸的单位范数约束。
\end{enumerate}
\end{learninggoals}

本章以前一章的连续嵌入为起点，但“连续”与“紧”所回答的问题不同。连续嵌入对每个函数给出同一常数的估计；紧嵌入则使任意有界列都能抽出强收敛子列。有限维空间中有界列自然具有这一性质，无限维函数空间却存在振荡、逃逸和集中，必须逐一控制。

阅读\secref{b1:sec:2501}时可以暂时忘记弱导数，只关心怎样用有限个平均值逼近整族函数。\secref{b1:sec:2502}再将这个判据用于 Sobolev 空间，所需延拓已在第22章建立。最后一节接回第9章的自反性与弱下半连续性：先问近似问题给了哪种有界估计，再问通过极限究竟需要哪种收敛，而不是一开始就把所有收敛方式混为一谈。

平移判据可对照 Hanche-Olsen 与 Holden 的《The Kolmogorov--Riesz compactness theorem》\cite[Section 3]{HancheOlsenHolden2010Compactness}阅读；本章把其有限维投影路线中的网格平均误差和远处截断分别展开。Sobolev 紧嵌入及变分方法的系统伴读可选 Adams 与 Fournier 的《Sobolev Spaces》\cite{Adams2003Sobolev}和 Brezis 的《Functional Analysis, Sobolev Spaces and Partial Differential Equations》\cite{Brezis2011Functional}。本章的证明始终从前面已经建立的工具出发，不要求读者预先掌握第二卷的偏微分方程理论。

% !TeX root = ../../Book1.tex
\section{\texorpdfstring{\(L^p\)}{Lᵖ} 中的平移判据}
\label{b1:sec:2501}

范数有界只能把函数限制在一个大球内，不能阻止它们向无穷远移动，也不能阻止固定区域内越来越快的振荡。强紧性需要同时限制这些变化。本节用远处的积分尾部和小平移误差刻画它，并把证明落实为有限个网格平均值的逼近；暂不要求函数具有弱导数。

Hanche-Olsen 与 Holden 的《The Kolmogorov--Riesz compactness theorem》\cite[Section 3, Theorem 5]{HancheOlsenHolden2010Compactness}采用有限维投影处理这一问题，并讨论不同文献中的定理命名。本节固定 \(d\geq1\)、\(1\leq p<\infty\) 和 \(\mathbb K=\mathbb R\) 或 \(\mathbb C\)。所有紧性均指 \(L^p\) 范数拓扑下的紧性，不是弱紧性。

% 实读 HancheOlsenHolden2010Compactness 的作者公开副本：
% https://www.math.ntnu.no/conservation/2009/037.pdf ，第3节Theorem5及完整证明，内页3--4。
% 正式发表：Expositiones Mathematicae 28(4) (2010),385--394，
% DOI10.1016/j.exmath.2010.03.001；作者副本首页2009是稿本日期，不作发表年。
% 本节将有限投影拆成全空间网格平均与有限cell截断，保留差集体积因子2^d。
% 必要性用有限网；任意域的零延拓结论只在Lp层面，未借后节Sobolev紧嵌入。

\subsection{平移连续性}
\label{b1:subsec:250101}

沿用 \(\tau_hu(x)=u(x-h)\)。平移是 \(L^p(\mathbb R^d)\) 上的线性等距映射，并满足
\[
 \tau_h\tau_k=\tau_{h+k},\quad
 \|\tau_hu-\tau_hv\|_p=\|u-v\|_p.
\]
由\cref{b1:lem:lp-translation-continuity}，每个固定的 \(u\in L^p\) 都有
\(\|\tau_hu-u\|_p\to0\)。但这里允许控制 \(h\) 的尺度随 \(u\) 改变；若面对一整个函数族，就需要一个适用于所有成员的尺度。

为把量词写清楚，对非空族 \(\mathcal F\subseteq L^p(\mathbb R^d)\) 记
\[
 \omega_{\mathcal F}(\delta)
 =\sup_{u\in\mathcal F}\sup_{|h|\leq\delta}
       \|\tau_hu-u\|_p.
\]
于是所需条件为 \(\omega_{\mathcal F}(\delta)\to0\)：给定误差后，先选同一个 \(\delta\)，再允许任取 \(u\in\mathcal F\) 和 \(|h|\leq\delta\)。这比逐个函数的平移连续性更强。

这种一致性可以控制函数在每个小网格内被常数代替时的误差。固定 \(\delta>0\)，用半开立方体
\[
 Q_{\delta,k}=\delta\bigl(k+[0,1)^d\bigr),
 \quad k\in\mathbb Z^d
\]
划分全空间；这些集合两两不交，并覆盖每个点。定义
\[
 (P_\delta u)(x)=\frac1{\delta^d}\int_{Q_{\delta,k}}u(y)\,\mathrm dy
 \quad(x\in Q_{\delta,k}).
\]
有限测度上的 Hölder 不等式保证每个平均值有定义。即使 \(u\) 是复值函数，平均值仍按普通复积分定义，不取共轭。

\begin{lemma}\label{b1:lem:lp-cell-average}
算子 \(P_\delta\) 是 \(L^p(\mathbb R^d)\) 上的线性投影，且
\(\|P_\delta u\|_p\leq\|u\|_p\)。此外，
\begin{equation}\label{b1:eq:lp-cell-average-error}
 \|u-P_\delta u\|_p^p
 \leq\frac1{\delta^d}
       \int_{[-\delta,\delta]^d}\|\tau_hu-u\|_p^p\,\mathrm dh.
\end{equation}
因此
\[
 \sup_{u\in\mathcal F}\|u-P_\delta u\|_p
 \leq2^{d/p}\omega_{\mathcal F}(\sqrt d\,\delta).
\]
\end{lemma}
\begin{proof}
对每个网格 \(Q=Q_{\delta,k}\)，平均值的 Hölder 估计给出
\[
 \delta^d\left|\frac1{\delta^d}\int_Q u\right|^p
 \leq\int_Q|u|^p.
\]
对 \(k\) 求和，得到收缩性。平均一个已在每个网格上为常数的函数，不会改变它，故 \(P_\delta^2=P_\delta\)。

对 \(x\in Q\)，同一估计应用于 \(u(x)-u(y)\)，得
\[
 |u(x)-P_\delta u(x)|^p
 \leq\frac1{\delta^d}\int_Q|u(x)-u(y)|^p\,\mathrm dy.
\]
先在 \(x\) 上积分并对网格求和，再作换元 \(y=x+h\)。同一个网格中的两个点满足 \(h\in[-\delta,\delta]^d\)，所以
\[
 \begin{aligned}
 \|u-P_\delta u\|_p^p
 &\leq\frac1{\delta^d}
      \sum_k\int_{Q_{\delta,k}}\int_{Q_{\delta,k}}
             |u(x)-u(y)|^p\,\mathrm dy\,\mathrm dx\\
 &\leq\frac1{\delta^d}
      \sum_k\int_{Q_{\delta,k}}\int_{[-\delta,\delta]^d}
             |u(x)-u(x+h)|^p\,\mathrm dh\,\mathrm dx\\
 &=\frac1{\delta^d}
      \int_{[-\delta,\delta]^d}\|\tau_{-h}u-u\|_p^p\,\mathrm dh.
 \end{aligned}
\]
这些是非负积分，Tonelli 定理允许求和及换序。积分区域关于原点对称，令 \(h\) 换成 \(-h\) 即得\eqref{b1:eq:lp-cell-average-error}。最后，该区域体积为 \((2\delta)^d\)，其中每个向量满足 \(|h|\leq\sqrt d\,\delta\)，从而得到因子 \(2^{d/p}\)。
\end{proof}

差向量所在的立方体边长为 \(2\delta\)，不是 \(\delta\)。因此上式右端不是一个已经归一化为一的平均；这个体积差正是因子 \(2^d\) 的来源。

\subsection{Fréchet–Kolmogorov 定理}
\label{b1:subsec:250102}

先回顾证明中需要的度量事实。若对每个 \(\varepsilon>0\)，一个集合 \(A\) 都能被有限个半径为 \(\varepsilon\) 的球覆盖，则称 \(A\) 为\term[quanyoujie]{全有界}[totally bounded]集；这些球的中心构成一个有限 \(\varepsilon\)-网。

\begin{lemma}\label{b1:lem:complete-total-bounded}
在完备度量空间中，一个子集相对紧，当且仅当它全有界。
\end{lemma}
\begin{proof}
紧闭包的任意固定半径球覆盖都有有限子覆盖，故必要性成立。反过来，若 \(A\) 全有界，则闭包 \(K=\overline A\) 也全有界：先用半径为 \(\varepsilon/2\) 的有限网覆盖 \(A\)，把半径扩大到 \(\varepsilon\) 即覆盖闭包。

对 \(K\) 中任意序列，逐次使用半径 \(2^{-j}\) 的有限覆盖，保留落在同一球中的无限子列，再对角抽选，所得子列为 Cauchy 列。由于 \(K\) 闭且环境完备，它在 \(K\) 中收敛。

这也给出开覆盖意义下的紧性。对 \(K\) 的任意开覆盖，必有 \(\eta>0\)，使每个 \(x\in K\) 的相对球 \(B_K(x,\eta)\) 包含于某个覆盖成员。否则可选 \(x_j\)，使 \(B_K(x_j,1/j)\) 不包含于任何单个成员；取收敛子列及包含其极限的开集，就会在充分大时得到相反结论。有限网的中心也可选在 \(K\) 中：从每个与 \(K\) 相交的网球中取一点作新中心，至多把半径扩大一倍。因此可用中心在 \(K\) 内的有限 \(\eta/2\)-网覆盖 \(K\)，为每个中心选择一个容纳其 \(\eta\)-球的覆盖成员，便得到有限子覆盖。
\end{proof}

现在引入\term[frechet-kolmogorov-dingli]{Fréchet–Kolmogorov 定理}[Fréchet--Kolmogorov compactness theorem]。条件分别写成范数、尾部与平移的形式，便于在具体问题中使用。

\begin{theorem}[Fréchet–Kolmogorov]\label{b1:thm:frechet-kolmogorov-lp}
设 \(1\leq p<\infty\)，\(\mathcal F\subseteq L^p(\mathbb R^d;\mathbb K)\) 非空。则 \(\mathcal F\) 相对紧，当且仅当下列条件都成立。
\begin{enumerate}
\item\label{b1:item:fk-bounded}
范数一致有界：
\[
 \sup_{u\in\mathcal F}\|u\|_p<\infty.
\]
\item\label{b1:item:fk-tail}
远处尾部一致趋零：
\[
 \lim_{R\to\infty}\sup_{u\in\mathcal F}
       \int_{|x|>R}|u(x)|^p\,\mathrm dx=0.
\]
\item\label{b1:item:fk-translation}
小平移误差一致趋零：
\[
 \lim_{\delta\downarrow0}\omega_{\mathcal F}(\delta)=0.
\]
\end{enumerate}
等价地，\(\mathcal F\) 中每个序列都有一个在 \(L^p\) 范数中收敛的子列，极限允许在 \(\overline{\mathcal F}\) 中。
\end{theorem}
\begin{proof}
先证明必要性。相对紧性给出任意尺度的有限网。一个有限 \(1\)-网已经说明范数有界，因为每个 \(u\) 的范数不超过其所靠近的网点范数加一。

固定 \(\varepsilon>0\)，取有限 \(\varepsilon\)-网
\(g_1,\ldots,g_N\)。每个 \(g_j\in L^p\) 都满足
\(\|\mathbf1_{\{|x|>R\}}g_j\|_p\to0\)。网点有限，可以选同一个 \(R\)，使这些尾部范数都小于 \(\varepsilon\)。给定 \(u\in\mathcal F\)，选取
\(\|u-g_j\|_p<\varepsilon\) 的网点，则
\[
 \|\mathbf1_{\{|x|>R\}}u\|_p
 \leq\|u-g_j\|_p+\|\mathbf1_{\{|x|>R\}}g_j\|_p<2\varepsilon.
\]
由此得到尾部条件。

同样，由每个网点的平移连续性，可选同一个 \(\delta>0\)，使
\(\|\tau_hg_j-g_j\|_p<\varepsilon\) 对全部 \(j\) 及
\(|h|\leq\delta\) 成立。平移等距性给出
\[
 \|\tau_hu-u\|_p
 \leq2\|u-g_j\|_p+\|\tau_hg_j-g_j\|_p<3\varepsilon.
\]
这就把有限多个函数的连续性变成了整族的一致连续性。

再证充分性。给定 \(\varepsilon>0\)，由平移条件选 \(\delta>0\)，使
\[
 \sup_{u\in\mathcal F}\|u-P_\delta u\|_p<\varepsilon/3.
\]
这是网格平均引理的直接应用，选取时保留其中的
\(2^{d/p}\) 及 \(\sqrt d\)。再由尾部条件选 \(R\)，使
\[
 \sup_{u\in\mathcal F}\|\mathbf1_{\{|x|>R\}}u\|_p<\varepsilon/3.
\]
保留所有与 \(B(0,R)\) 相交的网格，指标集合记为 \(J_{\delta,R}\)，它是有限集；令
\[
 S_{\delta,R}=\bigcup_{k\in J_{\delta,R}}Q_{\delta,k},
 \quad A_{\delta,R}u=\mathbf1_{S_{\delta,R}}P_\delta u.
\]
未保留的网格全部位于球外，至多涉及一个零测球面。逐网格使用平均值收缩估计，得到
\begin{equation}\label{b1:eq:cell-projection-tail}
 \begin{aligned}
 \|P_\delta u-A_{\delta,R}u\|_p^p
 &=\sum_{k\notin J_{\delta,R}}
       \delta^d\left|\frac1{\delta^d}\int_{Q_{\delta,k}}u\right|^p\\
 &\leq\sum_{k\notin J_{\delta,R}}\int_{Q_{\delta,k}}|u|^p
 \leq\int_{|x|>R}|u|^p.
 \end{aligned}
\end{equation}
因此
\(\|u-A_{\delta,R}u\|_p<2\varepsilon/3\)，对整族一致成立。

所有 \(A_{\delta,R}u\) 落在有限维空间
\[
 V_{\delta,R}
 =\operatorname{span}\{\,\mathbf1_{Q_{\delta,k}}:k\in J_{\delta,R}\,\},
\]
且范数一致有界。具体地，若 \(v=\sum_{k\in J_{\delta,R}}a_k\mathbf1_{Q_{\delta,k}}\)，则
\[
 \|v\|_p^p=\delta^d\sum_{k\in J_{\delta,R}}|a_k|^p.
\]
于是每个系数都落在固定的有界区间或有界复平面方框中。将这些有限个坐标分别用足够细的有限网格逼近，即得
\(A_{\delta,R}\mathcal F\) 的有限 \(\varepsilon/3\)-网；每个坐标误差不超过 \(\eta\) 时，总误差不超过
\(\delta^{d/p}|J_{\delta,R}|^{1/p}\eta\)。

三角不等式说明同一组网点构成 \(\mathcal F\) 的有限
\(\varepsilon\)-网。因此 \(\mathcal F\) 全有界。由\cref{b1:thm:lp-complete}及前一引理，闭包紧。引理证明中的抽选过程同时给出所述子列结论；反过来，若不存在某个尺度的有限网，就可逐个选出彼此距离至少为该尺度的序列，它不可能有收敛子列。
\end{proof}

证明始终先对原函数作全空间网格平均，然后截去平均函数的远处部分。\eqref{b1:eq:cell-projection-tail}控制了第二步；没有把“先截断函数、再平移”产生的边界差当成零。范数一致有界用于保证保留下来的有限个平均值有界，尾部与平移条件则决定需要保留多少网格、网格应当多细。

\subsection{尾部一致控制与支集}
\label{b1:subsec:250103}

若所有函数都在同一个有界集外几乎处处为零，尾部条件自动满足。反过来，尾部条件不要求严格的共同紧支集。例如，一个已经在 \(L^p\) 中收敛的序列连同极限是紧集，定理保证其尾部一致小，而各成员的支集完全可以无界。

固定支集不能排除振荡。令 \(Q=(0,1)^d\)，定义
\[
 u_n(x)=\mathbf1_Q(x)(-1)^{\lfloor2^n x_1\rfloor}.
\]
这些函数的 \(L^p\) 范数均为一，且在同一个紧集外为零。当 \(m>n\) 时，在每个较粗的二进小区间内，两种符号相同与相反的部分各占一半，故
\[
 \|u_m-u_n\|_p^p=2^{p-1}.
\]
所以任何子列都不是 Cauchy 列。还可直接看出平移条件失败：取
\(h_n=2^{-n}e_1\)，在
\((2^{-n},1)\times(0,1)^{d-1}\) 上，\(u_n(x-h_n)=-u_n(x)\)，从而
\[
 \|\tau_{h_n}u_n-u_n\|_p^p
 \geq2^p(1-2^{-n}).
\]
每个固定的 \(u_n\) 仍有自己的平移连续性，但这些尺度不能统一。

另一种失败是向远处移动。取 \(g=\mathbf1_Q\)，令
\(v_n=\tau_{3ne_1}g\)。由于平移可交换，
\[
 \|\tau_hv_n-v_n\|_p=\|\tau_hg-g\|_p,
\]
小平移条件对整列一致成立。然而对每个固定 \(R\)，充分远的
\(v_n\) 把全部 \(L^p\) 质量放在 \(B(0,R)\) 外，尾部条件失败。其支集两两不交，故不同项之间的距离为 \(2^{1/p}\)。这与前面的振荡例子是不同的障碍。

现在把判据用于一般定义域。这里只需要 Lebesgue 可测性，不要求边界光滑。

\begin{corollary}\label{b1:cor:lp-domain-zero-extension-compactness}
设 \(\Omega\subseteq\mathbb R^d\) 为 Lebesgue 可测集，且
\(\mathcal F\subseteq L^p(\Omega;\mathbb K)\)。把每个 \(u\) 在域外补零，记为 \(E_0u\)。则 \(\mathcal F\) 在 \(L^p(\Omega)\) 中相对紧，当且仅当它范数一致有界，且
\[
 \lim_{R\to\infty}\sup_{u\in\mathcal F}
      \int_{\Omega\cap\{|x|>R\}}|u|^p=0,\quad
 \lim_{\delta\downarrow0}\sup_{u\in\mathcal F}
      \sup_{|h|\leq\delta}\|\tau_hE_0u-E_0u\|_{L^p(\mathbb R^d)}=0.
\]
若 \(\Omega\) 有界，第一个极限条件自动成立。空族按平凡情形处理。
\end{corollary}
\begin{proof}
补零是 \(L^p(\Omega)\) 到 \(L^p(\mathbb R^d)\) 的线性等距嵌入，其像为在 \(\Omega^c\) 上几乎处处为零的函数所成的闭子空间。闭性可直接检查：若 \(E_0u_j\to f\) 于全空间 \(L^p\)，则
\[
 \|\mathbf1_{\Omega^c}f\|_p
 \leq\|f-E_0u_j\|_p\longrightarrow0.
\]
因此 \(\mathcal F\) 相对紧与 \(E_0\mathcal F\) 相对紧等价。把前一定理用于补零后的族，三个条件正好成为所列条件。
\end{proof}

其中的平移误差必须在整个空间计算。若只在两份定义域的交集上比较函数，就会漏掉跨越边界的部分。事实上，记
\(\Omega+h=\{x+h:x\in\Omega\}\)，分割积分区域可得
\[
 \begin{aligned}
 \|\tau_hE_0u-E_0u\|_p^p
 &=\int_{\Omega\cap(\Omega+h)}
         |u(x-h)-u(x)|^p\,\mathrm dx\\
 &\quad+\int_{\Omega\setminus(\Omega+h)}|u(x)|^p\,\mathrm dx
       +\int_{\Omega\setminus(\Omega-h)}|u(x)|^p\,\mathrm dx.
 \end{aligned}
\]
最后一项由 \(x\in(\Omega+h)\setminus\Omega\) 的积分作平移换元得到。这个恒等式既记录了域内变化，也记录了靠近边界的质量；此处从未假设弱导数可以补零。后续处理 Sobolev 函数时，还须另用延拓定理或零边界条件建立所需的平移估计。

有限 \(p\) 的限制不能直接删去。\secref{b1:sec:1304}已经指出，
\(u=\mathbf1_{(0,1)}\) 满足
\(\|\tau_hu-u\|_\infty=1\)（\(0<|h|<1\)）。因此连单元素紧集
\(\{u\}\) 都不满足上述平移条件的无穷范数版本。网格平均的积分误差证明也没有提供对应的上确界逼近，不能把本定理机械改成 \(p=\infty\)。

% !TeX root = ../../Book1.tex
\section{Rellich–Kondrachov 定理}
\label{b1:sec:2502}

第24章的嵌入估计说明，一族函数的弱导数若有统一积分界，函数本身就有更强的可积性或连续性。但连续嵌入只给出范数估计，尚未保证有收敛子列。本节再用梯度控制小平移，从而把 Sobolev 有界性接到上一节的紧性判据。区域有界提供尾部控制，边界延拓则使平移能够在全空间进行。

两项限制会贯穿证明。第一，一般区域上的函数不能直接补零后求弱导数，必须使用第22章建立的延拓。第二，从连续嵌入得到紧嵌入时，目标指数通常要留出严格余量；临界指数处的缩放会保留一部分范数，而使支集越来越小。

\subsection{有界区域}
\label{b1:subsec:250201}

先明确梯度所提供的统一模。以下有限 \(p\) 的估计已用于第23章的整数阶到分数阶嵌入；把它单列出来，便于识别紧性证明的输入。

\begin{lemma}\label{b1:lem:sobolev-translation-gradient}
若 \(1\le p<\infty\)、\(u\in W^{1,p}(\mathbb R^d)\)，则
\begin{equation}\label{b1:eq:sobolev-translation-gradient}
 \|\tau_hu-u\|_{L^p(\mathbb R^d)}
 \le |h|\cdot\bigl\|\,|\nabla u|\,\bigr\|_{L^p(\mathbb R^d)}
 \quad(h\in\mathbb R^d).
\end{equation}
\end{lemma}
\begin{proof}
先设 \(u\in C_c^\infty(\mathbb R^d)\)。沿线段积分得
\[
 u(x-h)-u(x)=-\int_0^1\nabla u(x-th)\cdot h\,\mathrm dt.
\]
对 \(t\) 的概率积分使用 Jensen 不等式，再对 \(x\) 积分，得到
\[
 \int_{\mathbb R^d}|u(x-h)-u(x)|^p\,\mathrm dx
 \le |h|^p\int_0^1\int_{\mathbb R^d}|\nabla u(x-th)|^p
                \,\mathrm dx\,\mathrm dt
 =|h|^p\int_{\mathbb R^d}|\nabla u|^p.
\]
由\cref{b1:thm:whole-space-sobolev-density}取光滑逼近；平移等距，梯度的欧氏 \(L^p\) 范数又受标准 Sobolev 范数控制，故两端都能通过极限。
\end{proof}

\begin{proposition}\label{b1:prop:rellich-base-lp}
设 \(\Omega\subset\mathbb R^d\) 为有界 Lipschitz 区域，\(1\le p<\infty\)。则 \(W^{1,p}(\Omega)\) 的任意有界子集在 \(L^p(\Omega)\) 中相对紧。
\end{proposition}
\begin{proof}
由\cref{b1:thm:lipschitz-domain-extension}，存在固定的有界线性延拓算子 \(E\)，其像中所有函数都在一个与函数无关的紧集 \(K\) 外为零。若 \(\mathcal F\) 在 \(W^{1,p}(\Omega)\) 中有界，则 \(E\mathcal F\) 在全空间 \(W^{1,p}\) 中有界，因而在 \(L^p\) 中有界；共同支集保证远处尾部为零；上一引理又给出
\[
 \sup_{u\in\mathcal F}\|\tau_hEu-Eu\|_p\le C|h|.
\]
Fréchet–Kolmogorov 定理的三个条件都满足。因此任取 \(u_j\in\mathcal F\)，可抽取子列使 \(Eu_j\) 在全空间 \(L^p\) 中收敛。限制到 \(\Omega\) 不增加 \(L^p\) 距离，就得到所需收敛子列。
\end{proof}

证明不要求 \(\Omega\) 连通，因为所用的一阶延拓不要求连通。它也不要求事先知道序列有弱极限，因而包括 \(p=1\)。不过，把“有界 Lipschitz 区域”仅改成“有界开集”，对一般 \(W^{1,p}\) 就不成立；第\ref{b1:ex:25-infinitely-many-components}题用越来越小的连通分支给出反例。

\subsection{次临界嵌入}
\label{b1:subsec:250202}

若赋范空间之间的包含映射连续，并把有界集映为相对紧集，就称这个嵌入为\term[jinqianru]{紧嵌入}[compact embedding]。下面的结果通常称为\term[rellich-kondrachov-dingli]{Rellich–Kondrachov 定理}[Rellich--Kondrachov theorem]。这里分别列出积分范数与连续函数范数的结论。

\begin{theorem}[Rellich–Kondrachov，有限目标指数]\label{b1:thm:rellich-subcritical}
设 \(\Omega\subset\mathbb R^d\) 为有界 Lipschitz 区域，\(1\le p<\infty\)。下列嵌入是紧的：
\begin{enumerate}
\item 若 \(p<d\)，则 \(W^{1,p}(\Omega)\subset L^q(\Omega)\)，其中 \(1\le q<dp/(d-p)\)。
\item 若 \(p=d\)，则 \(W^{1,p}(\Omega)\subset L^q(\Omega)\)，其中 \(1\le q<\infty\)。
\item 若 \(p>d\)，则 \(W^{1,p}(\Omega)\subset L^q(\Omega)\)，其中 \(1\le q<\infty\)；实际上连续代表还满足下面的更强结论。
\end{enumerate}
\end{theorem}
\begin{proof}
取 Sobolev 有界列。前一命题给出在 \(L^p\) 中收敛的子列。当 \(q=p\) 时，这正是前一命题的结论。以下只需处理 \(q\ne p\)。若 \(q<p\)，有限测度上的 Hölder 不等式
\[
 \|v\|_q\le m_d(\Omega)^{1/q-1/p}\|v\|_p
\]
已经给出 \(L^q\) 收敛。若存在 \(r>q>p\)，使原列的 \(L^r\) 范数一致有界，则对任意两项之差使用第10章的插值估计，得到
\[
 \|u_j-u_k\|_q
 \le \|u_j-u_k\|_p^\theta
       \cdot\|u_j-u_k\|_r^{1-\theta},
 \qquad
 \frac1q=\frac\theta p+\frac{1-\theta}r.
\]
于是子列在 \(L^q\) 中也是 Cauchy 列。由完备性它有 \(L^q\) 极限；有限测度上两种范数收敛都蕴含依测度收敛，故极限相同。

当 \(p<d\) 时，由第24章的域上 Sobolev 嵌入可取 \(r=p^*=dp/(d-p)\)。当 \(p=d\ge2\) 时，\secref{b1:sec:2405}已经证明任意有限指数的连续嵌入，任选有限的 \(r>q\) 即可。若不读该选读节，也可先在有限测度的 \(\Omega\) 上降到 \(p_0<d\)，选 \(p_0\) 足够接近 \(d\) 使 \(dp_0/(d-p_0)>q\)，再用次临界嵌入获得同一 \(L^r\) 界。

剩下 \(d=p=1\)。利用一阶延拓取得共同有界支集的 \(W^{1,1}(\mathbb R)\) 函数。取一个包含这些支集的有限区间；由主线\secref{b1:sec:1604}中的绝对连续代表与微积分基本定理（\cref{b1:thm:ac-fundamental-calculus}），其代表满足
\[
 |Eu(x)|\le\int_{\mathbb R}|(Eu)'(t)|\,\mathrm dt.
\]
这是从支集左侧的零值沿区间积分得到的。因此有统一 \(L^\infty\) 界，在有限测度的 \(\Omega\) 上也有任意有限 \(L^r\) 界，同样完成插值。

若 \(p>d\)，Morrey 不等式给出连续代表的统一上确界界，再取有限 \(r>q\) 即可。上述各步也同时给出包含映射的连续性。
\end{proof}

这个证明的紧性来源是 \(L^p\) 中的平移控制；更高指数来自第24章已经建立的连续估计。插值中的 \(\theta>0\) 很重要：正是这一幂次把已知趋零的 \(L^p\) 差保留下来。到 \(q=p^*\) 时，不能再由此获得趋零因子。

\begin{theorem}[Rellich–Kondrachov，连续代表]\label{b1:thm:rellich-holder}
设 \(\Omega\) 为有界 Lipschitz 区域。
若 \(d<p<\infty\)，记 \(\alpha=1-d/p\)；若 \(p=\infty\)，记 \(\alpha=1\)。
则连续代表给出的映射
\[
 W^{1,p}(\Omega)\longrightarrow C(\overline\Omega)
 \quad\text{及}\quad
 W^{1,p}(\Omega)\longrightarrow C^{0,\beta}(\overline\Omega)
 \quad(0<\beta<\alpha)
\]
都是紧嵌入。其中 \(C(\overline\Omega)\) 取上确界范数。
\end{theorem}
\begin{proof}
由延拓及 Morrey 估计，有限 \(p>d\) 时，有界列的连续代表满足
\[
 \sup_j\|u_j\|_{C^{0,\alpha}(\overline\Omega)}\le M.
\]
若 \(p=\infty\)，第21章的 Lipschitz 代表与第22章的全空间延拓给出同样的结论，指数取 \(1\)。这时并未把有限 \(p\) 的平移判据推广到无穷范数。

下面直接证明所需的一致收敛子列。紧集 \(\overline\Omega\) 在每个尺度 \(1/n\) 都有有限网；把这些网点排成一个可数集。每个网点处的数列有界，对网点依次抽取子列并取对角列，便使所有网点处的值都收敛。给定 \(\varepsilon>0\)，先取一个足够细的有限网，使 \(2M|x-y|^\alpha<\varepsilon/2\) 对相邻点与网点成立；再把这有限个网点处的两项差控制在 \(\varepsilon/2\) 内。三角不等式说明对角列在整个 \(\overline\Omega\) 上一致 Cauchy。逐点取极限并使用这一统一控制，所得极限连续，子列一致收敛。

还须增强到较低阶 Hölder 范数。对任意 \(v\in C^{0,\alpha}(\overline\Omega)\)，令 \(A=\|v\|_\infty\)、\(B=[v]_{0,\alpha}\)。当 \(A,B>0\) 时，对不同的两点及 \(r=|x-y|\)，
\[
 |v(x)-v(y)|\le\min\{2A,Br^\alpha\}.
\]
在 \(r=(2A/B)^{1/\alpha}\) 两侧分别使用两个估计，得到
\begin{equation}\label{b1:eq:holder-interpolation-compact}
 [v]_{0,\beta}\le
 (2A)^{1-\beta/\alpha}B^{\beta/\alpha}.
\end{equation}
若 \(A=0\) 或 \(B=0\)，结论直接成立。子列的一致极限仍有 Hölder 半范数不超过 \(M\)，因为每个固定两点的估计都能通过极限。对 \(v=u_j-u\) 使用\eqref{b1:eq:holder-interpolation-compact}，右端的 \(A\) 趋零而 \(B\le2M\)，于是得到 \(C^{0,\beta}\) 收敛。
\end{proof}

这里需要 \(\beta<\alpha\)，与积分情形的严格余量相对应。第\ref{b1:ex:25-critical-mass-concentration}题和第\ref{b1:ex:25-holder-endpoint}题分别计算临界积分指数与临界 Hölder 指数的集中例子：有界性仍在，但紧性丧失。另一方面，\(W^{1,1}\) 在区间上虽由绝对连续代表连续嵌入 \(L^\infty\)，其有界列并没有统一的连续模，因此本定理没有把 \(p=d=1\) 列为连续函数空间中的紧嵌入。

\subsection{边界条件}
\label{b1:subsec:250203}

区域边界的规则性与函数自身的边界条件，可以分别提供进入全空间所需的通道。对零边界闭包，补零已由第22章证明能够保留弱导数，因此不必再次要求 Lipschitz 边界。

\begin{corollary}\label{b1:cor:rellich-zero-boundary}
设 \(\Omega\subset\mathbb R^d\) 为任意有界开集。若 \(1\le p<\infty\)，则下列嵌入是紧的：
\begin{enumerate}
\item 当 \(p<d\) 时，\(W_0^{1,p}(\Omega)\subset L^q(\Omega)\)，其中 \(1\le q<dp/(d-p)\)；
\item 当 \(p=d\) 时，\(W_0^{1,p}(\Omega)\subset L^q(\Omega)\)，其中 \(1\le q<\infty\)；
\item 当 \(p>d\) 时，\(W_0^{1,p}(\Omega)\subset L^q(\Omega)\)，其中 \(1\le q<\infty\)。
\end{enumerate}
若 \(p=\infty\)，则 \(W_0^{1,\infty}(\Omega)\subset L^q(\Omega)\) 对每个 \(1\le q<\infty\) 也是紧嵌入。进一步，当 \(d<p<\infty\) 时令 \(\alpha=1-d/p\)，当 \(p=\infty\) 时令 \(\alpha=1\)；补零后的连续代表限制到 \(\overline\Omega\)，给出
\[
 W_0^{1,p}(\Omega)\longrightarrow C(\overline\Omega),
 \qquad
 W_0^{1,p}(\Omega)\longrightarrow C^{0,\beta}(\overline\Omega)
 \quad(0<\beta<\alpha)
\]
这两类紧嵌入。
\end{corollary}
\begin{proof}
有限 \(p\) 时，\cref{b1:prop:w0-zero-extension}把有界列等距送入 \(W^{1,p}(\mathbb R^d)\)，共同支集包含于紧集 \(\overline\Omega\)。梯度平移估计和 Fréchet–Kolmogorov 定理直接给出全空间 \(L^p\) 收敛子列。更高积分指数的界，取第24章的零边界版本；临界 \(p=d\ge2\) 也可先在包含 \(\overline\Omega\) 的球上降低指数，再用该球的连续嵌入。其余插值步骤完全相同。

当 \(p>d\) 时，对全空间补零函数使用 Morrey 估计；\(p=\infty\) 时使用全空间 Lipschitz 代表。于是有限网抽列及 Hölder 插值的证明都直接用于紧集 \(\overline\Omega\)。所得一致收敛又在有限测度的 \(\Omega\) 上蕴含每个有限 \(L^q\) 范数中的收敛。
\end{proof}

上述连续代表确实在每个边界点为零，而不只是在域外几乎处处为零。由 \(W_0^{1,p}\) 的闭包定义，可用 \(C_c^\infty(\Omega)\) 函数在 Sobolev 范数下逼近。其补零函数在全空间有同一范数收敛；对 \(p>d\) 的 Morrey 估计或 \(p=\infty\) 的上确界控制，使这些补零函数一致趋于连续代表。每个逼近函数在 \(\partial\Omega\) 都为零，极限也为零。这个理由包括边界中可能存在的裂缝，不能只凭“域外几乎处处为零”略过。

在有界 Lipschitz 区域上，有限 \(p\) 的零迹刻画把这个推论与第23章接起来。固定边界迹也不会破坏紧性：若 \(g\in W^{1,p}(\Omega)\) 固定，而 \(u_j-g\in W_0^{1,p}(\Omega)\) 且 \(u_j\) 有界，就对 \(u_j-g\) 使用推论，再把 \(g\) 加回去。在任意有界开集上，同样可以用这个闭仿射条件表述边界约束；不必先给粗糙边界定义迹。

内部紧性则完全不需要全域边界条件。若 \(1\le p<\infty\)，且 \((u_j)\) 在 \(W^{1,p}_{\mathrm{loc}}(\Omega)\) 中局部有界，则可抽取共同子列，使它在每个 \(U\Subset\Omega\) 上强收敛于 \(L^p(U)\)。利用同一子列及连续嵌入与插值，还可得到：当 \(p<d\) 时，对每个 \(1\le q<dp/(d-p)\) 强收敛于 \(L^q(U)\)；当 \(p=d\) 时，对每个有限 \(q\) 成立；当 \(p>d\) 时，同样对每个有限 \(q\) 成立。若 \(d<p<\infty\)，记 \(\alpha=1-d/p\)；另若 \((u_j)\) 在 \(W^{1,\infty}_{\mathrm{loc}}(\Omega)\) 中局部有界，则记 \(\alpha=1\)。在这两种情形中可取连续代表，使相应的共同子列在每个 \(U\Subset\Omega\) 上局部一致收敛，并在每个 \(C^{0,\beta}(\overline U)\)、\(0<\beta<\alpha\) 中收敛。

证明这些局部结论时，固定 \(\overline U\Subset\Omega\)，取等于一于 \(\overline U\) 邻域的光滑紧支集截断。乘积补零后形成一个具有共同紧支集的全空间 Sobolev 有界族，故可使用前述相应结论；再对可数内部穷尽作对角抽选。所得结论不保证原区域的全局范数收敛，因为质量仍可能靠近边界或移向无穷远。

\subsection{高阶版本}
\label{b1:subsec:250204}

高阶紧性首先表现为损失一阶导数后的强收敛。为了不预借尚未证明的高阶边界延拓，先把所需条件写在命题中。

\begin{proposition}\label{b1:prop:rellich-higher-order}
设 \(k\ge2\)、\(1\le p<\infty\)，\(\Omega\) 为有界开集。若存在有界线性延拓
\[
 E_k:W^{k,p}(\Omega)\longrightarrow W^{k,p}(\mathbb R^d),
 \qquad (E_ku)|_\Omega=u,
\]
则 \(W^{k,p}(\Omega)\subset W^{k-1,p}(\Omega)\) 是紧嵌入。
无论是否存在这个延拓，令
\[
 W_0^{k,p}(\Omega)
 =\overline{C_c^\infty(\Omega)}^{\,W^{k,p}(\Omega)},
\]
则 \(W_0^{k,p}(\Omega)\subset W^{k-1,p}(\Omega)\) 总是紧嵌入。
\end{proposition}
\begin{proof}
先处理延拓情形。取等于一于 \(\overline\Omega\) 邻域的固定光滑紧支集函数 \(\chi\)。乘子估计说明 \(u\mapsto\chi E_ku\) 有界进入全空间 \(W^{k,p}\)，而且所有像具有共同紧支集。对 \(|\gamma|\le k-1\)，函数 \(\partial^\gamma(\chi E_ku_j)\) 在全空间 \(W^{1,p}\) 中有界，故由梯度平移界和紧性判据可抽出 \(L^p\) 收敛子列。这样的 \(\gamma\) 只有有限多个，逐次抽取可使它们同时收敛。

限制到 \(\Omega\) 后，所得子列在 \(W^{k-1,p}(\Omega)\) 范数中 Cauchy，因为这个范数由刚才有限多个分量的 \(p\) 次和构成。完备性给出该空间中的极限；或者逐项对测试函数通过积分极限，也能确认低阶导数的相容关系。

再看零边界闭包。对 \(u\in W_0^{k,p}(\Omega)\)，取 \(C_c^\infty(\Omega)\) 逼近列。每一项及其各阶导数补零后，在全空间的 \(W^{k,p}\) 距离与原域距离相同，因而为 Cauchy 列。由全空间完备性得到极限；其零阶分量是 \(E_0u\)，所有导数均为原弱导数的补零。因此补零是所需的等距延拓，并具有共同有界支集，前面的有限分量抽列再次适用。
\end{proof}

若关心更高的积分指数，可在上述证明中对每个低阶导数使用第24章的高阶连续嵌入，再与已获得的强 \(L^p\) 收敛插值。具体地，若 \(kp<d\) 且命题中的有界延拓 \(E_k\) 存在，则
\[
 W^{k,p}(\Omega)\subset L^q(\Omega),
 \qquad 1\le q<\frac{dp}{d-kp}
\]
为紧嵌入；对任意有界开集，不要求存在 \(E_k\)，也有
\[
 W_0^{k,p}(\Omega)\subset L^q(\Omega),
 \qquad 1\le q<\frac{dp}{d-kp}
\]
为紧嵌入。右端的严格不等式仍不能省去。

同样，设 \(k-d/p=m+\alpha>0\)，其中 \(m\) 为非负整数、\(0<\alpha<1\)。若有上述延拓 \(E_k\)，则
\[
 W^{k,p}(\Omega)\longrightarrow C^{m,\beta}(\overline\Omega)
 \quad(0<\beta<\alpha)
\]
为紧嵌入；对任意有界开集，相应的
\[
 W_0^{k,p}(\Omega)\longrightarrow C^{m,\beta}(\overline\Omega)
 \quad(0<\beta<\alpha)
\]
也为紧嵌入。第24章给出全部不超过 \(m\) 阶导数的统一连续界及 \(m\) 阶导数的统一 Hölder 界；对这有限多个分量使用有限网抽列和\eqref{b1:eq:holder-interpolation-compact}即可。极限分量的导数关系由内部线段积分通过一致极限确认，与第24章相同。若 \(k-d/p\) 为正整数，则须在该章已证明的低于 \(1\) 的 Hölder 指数中另留严格余量，并分别保留上述延拓或零边界条件。

这些表述中的高阶延拓假设没有因为使用了第24章而消失。对于任意开集，可以始终取得内部版本；对于零边界闭包，可以始终使用补零；对于一般高阶 Sobolev 函数的全边界版本，还必须给出对应的延拓。把不同通道分开，才能看清紧性究竟由哪些已建立的条件保证。

% !TeX root = ../../Book1.tex
\section{弱收敛方法}
\label{b1:sec:2503}

紧嵌入使有界序列在较弱的范数下具有收敛子列，Sobolev 空间的自反性则在原空间中提供弱收敛子列。这两种结论经常一起使用，但承担不同的任务：弱收敛保留导数的积分恒等式，紧嵌入加强函数本身的收敛。至于某个非线性能量能否通过极限，还要看它的凸性、增长条件以及所需的收敛方式。

除最后的实值变分问题外，本节允许实值或复值函数。主要指数范围为 \(1<p<\infty\)，并记 \(p'=p/(p-1)\)。
% 本节证明直接接本书第9章弱闭性、凸泛函弱下半连续性，
% 第21章Sobolev分量结构及本章Rellich结论。
% Brezis2011Functional在来源台账中只有出版社信息与目录证据；
% 文末仅作系统伴读推荐，不将未取得的全文列为某项证明的实读来源。

\subsection{有界序列}
\label{b1:subsec:250301}

设 \(\Omega\subset\mathbb R^d\) 为开集，并已从一个近似问题中取得估计
\[
 \sup_j\left(
   \|u_j\|_{L^p(\Omega)}^p+
   \sum_{i=1}^d\|\partial_i u_j\|_{L^p(\Omega)}^p
 \right)<\infty.
\]
由\cref{b1:prop:sobolev-separable-reflexive}与\ref{b1:cor:sobolev-weak-subsequence}，可抽取子列，使
\[
 u_{j_\ell}\rightharpoonup u
 \quad\text{于 }W^{1,p}(\Omega).
\]
这一抽列不要求 \(\Omega\) 有界，也不要求边界光滑。它来自空间的自反性，而不是来自区域上的紧嵌入。

计算中也常先把函数和导数分别处理。由于分量只有有限个，对
\(u_j,\partial_1u_j,\ldots,\partial_du_j\) 逐次抽取子列，就能得到一个共同子列和函数 \(v_0,v_1,\ldots,v_d\in L^p(\Omega)\)，使各分量分别弱收敛到这些函数。此时还不能把 \(v_i\) 直接记成 \(\partial_i v_0\)：它们最初只是不同 \(L^p\) 空间中的极限，导数关系需要由测试函数确认。

在两个端点，上述自反性理由不再适用。当 \(p=1\) 时，有界导数的积分可能集中；\(p=\infty\) 时，利用 \(L^1\) 测试得到的弱-*收敛也不等于由整个连续对偶测试的弱收敛。前者的集中问题可回看\secref{b1:sec:1005}，后者的区别见\cref{b1:prop:sobolev-weak-components}后的说明。

\subsection{弱极限}
\label{b1:subsec:250302}

\begin{proposition}\label{b1:prop:sobolev-weak-limit-identification}
设 \(1<p<\infty\)、\(u_j\in W^{1,p}(\Omega)\)，并有
\[
 u_j\rightharpoonup v_0,\quad
 \partial_i u_j\rightharpoonup v_i
 \quad\text{于 }L^p(\Omega),\quad 1\le i\le d.
\]
则 \(v_0\in W^{1,p}(\Omega)\)、\(\partial_i v_0=v_i\)，且
\(u_j\rightharpoonup v_0\) 于 \(W^{1,p}(\Omega)\)。
若另有固定的 \(g\in W^{1,p}(\Omega)\) 及范数闭线性子空间
\(M\subset W^{1,p}(\Omega)\)，并且所有 \(u_j\in g+M\)，则 \(v_0\in g+M\)。
\end{proposition}
\begin{proof}
对任意 \(\varphi\in C_c^\infty(\Omega)\)，函数 \(\varphi,\partial_i\varphi\) 都属于 \(L^{p'}(\Omega)\)。将它们代入两个弱极限，得到
\[
 \int_\Omega v_0\partial_i\varphi\,\mathrm dx
 =\lim_{j\to\infty}\int_\Omega u_j\partial_i\varphi\,\mathrm dx
 =-\lim_{j\to\infty}\int_\Omega\partial_i u_j\varphi\,\mathrm dx
 =-\int_\Omega v_i\varphi\,\mathrm dx.
\]
这正是 \(v_i\) 为 \(v_0\) 的弱偏导数的定义。由\cref{b1:prop:sobolev-weak-components}，全部分量的弱收敛给出 Sobolev 空间中的弱收敛。

仿射集 \(g+M\) 范数闭且凸，故由\cref{b1:thm:convex-weak-closure}也是弱闭的。因此弱极限仍属于这个集合。
\end{proof}

取 \(M=W_0^{1,p}(\Omega)\)，便知边界约束
\[
 u_j-g\in W_0^{1,p}(\Omega)
\]
能够随弱极限保留。这里使用的是\cref{b1:def:sobolev-zero-boundary}中的范数闭包，不需要给任意粗糙边界上的函数指定逐点值。在有界 Lipschitz 区域上，若再调用迹定理，才可把这一约束解释为与 \(g\) 具有相同的迹。

导数的识别只用了线性积分关系。若近似方程中还出现
\(|\nabla u_j|^{p-2}\nabla u_j\) 等非线性表达式，不能仅凭
\(\nabla u_j\rightharpoonup\nabla u\) 就把它的弱极限写成对应的极限表达式。后面的变分论证将先比较能量，再从极小性质推导方程。

\subsection{强收敛的恢复}
\label{b1:subsec:250303}

已有弱极限时，紧性不只是再给出一条强收敛子列。它能确定像序列的全部极限，从而使整列收敛。第\ref{b1:ex:09-compact-improves}题中的抽列方法在此成为使用紧嵌入的基本步骤。

\begin{proposition}\label{b1:prop:compact-weak-to-strong}
设 \(X,Y\) 为赋范空间，\(T:X\rightarrow Y\) 为紧线性算子。若
\(x_j\rightharpoonup x\) 于 \(X\)，则 \(Tx_j\to Tx\) 于 \(Y\) 的范数。
\end{proposition}
\begin{proof}
弱收敛列有界。又因每个 \(y^*\in Y^*\) 与 \(T\) 的复合都属于 \(X^*\)，有
\(Tx_j\rightharpoonup Tx\)。若范数收敛不成立，可选 \(\varepsilon>0\) 及子列，使
\[
 \|Tx_{j_\ell}-Tx\|_Y\ge\varepsilon
 \quad\text{对所有 }\ell.
\]
紧性使这条子列还有范数收敛子列，设其极限为 \(y\)。范数收敛蕴含弱收敛，弱极限唯一，故 \(y=Tx\)，与上式矛盾。
\end{proof}

例如，设 \(\Omega\) 为有界 Lipschitz 区域，且
\(u_j\rightharpoonup u\) 于 \(W^{1,p}(\Omega)\)。由\cref{b1:thm:rellich-subcritical}，当 \(p<d\) 时，对每个 \(1\le q<p^*=dp/(d-p)\)，整列 \(u_j\to u\) 于 \(L^q(\Omega)\)；当 \(p=d\) 时，这对每个有限 \(q\) 成立。若 \(p>d\)，则\cref{b1:thm:rellich-holder}使连续代表在
\(C(\overline\Omega)\) 以及每个 \(0<\beta<1-d/p\) 对应的
\(C^{0,\beta}(\overline\Omega)\) 中整列收敛。这些结论没有使梯度自动强收敛。

要恢复原 Sobolev 范数的收敛，可以补充范数大小的信息。先证明所需的标量事实。

\begin{lemma}\label{b1:lem:lp-weak-norm-strong}
设 \(1<p<\infty\)。若 \(f_j\rightharpoonup f\) 于实或复 \(L^p(\mu)\)，且
\(\|f_j\|_p\to\|f\|_p\)，则 \(\|f_j-f\|_p\to0\)。
\end{lemma}
\begin{proof}
标量函数 \(a\mapsto|a|^p\) 严格凸：三角不等式的等号首先要求两个非零标量同向，而实函数 \(t\mapsto t^p\) 在非负半轴上的严格凸性又排除不同的长度。因此
\[
 D(a,b)=\frac{|a|^p+|b|^p}{2}
              -\left|\frac{a+b}{2}\right|^p
\]
非负，且仅在 \(a=b\) 时为零。固定 \(0<\varepsilon<1\)。在有限维紧集
\[
 \left\{\,(a,b):|a|^p+|b|^p=1,\ |a-b|\ge\varepsilon\,\right\}
\]
上，\(D\) 有正的最小值 \(\delta_\varepsilon\)。由齐次性，若
\(S=|a|^p+|b|^p\) 且 \(|a-b|\ge\varepsilon S^{1/p}\)，便有
\(D(a,b)\ge\delta_\varepsilon S\)。

另一方面，\((f_j+f)/2\rightharpoonup f\)，故范数弱下半连续性与三角不等式给出
\[
 \|f\|_p
 \le\liminf_j\left\|\frac{f_j+f}{2}\right\|_p
 \le\limsup_j\left\|\frac{f_j+f}{2}\right\|_p
 \le\|f\|_p.
\]
因此 \(\int D(f_j,f)\,\mathrm d\mu\to0\)。在
\(|f_j-f|<\varepsilon\bigl(|f_j|^p+|f|^p\bigr)^{1/p}\) 的部分，用该不等式控制差；在其余部分，使用刚才的亏损估计以及
\(|a-b|^p\le2^{p-1}S\)。积分后得到
\[
 \|f_j-f\|_p^p
 \le\varepsilon^p\bigl(\|f_j\|_p^p+\|f\|_p^p\bigr)
   +\frac{2^{p-1}}{\delta_\varepsilon}
       \int D(f_j,f)\,\mathrm d\mu.
\]
先让 \(j\to\infty\)，再让 \(\varepsilon\downarrow0\)，即得结论。
\end{proof}

\begin{proposition}\label{b1:prop:sobolev-weak-norm-strong}
设 \(1<p<\infty\)，且 \(u_j\rightharpoonup u\) 于 \(W^{1,p}(\Omega)\)。若本书规定的标准范数满足
\[
 \|u_j\|_{W^{1,p}(\Omega)}
 \longrightarrow\|u\|_{W^{1,p}(\Omega)},
\]
则 \(u_j\to u\) 于 \(W^{1,p}(\Omega)\)。
\end{proposition}
\begin{proof}
在 \(Z=\Omega\times\{0,\ldots,d\}\) 上取 Lebesgue 测度与有限计数测度的乘积，令
\[
 F_j(x,0)=u_j(x),\quad F_j(x,i)=\partial_i u_j(x)
 \quad(1\le i\le d),
\]
并由 \(u\) 定义 \(F\)。每个 \(L^{p'}(Z)\) 测试函数只有有限多个分量，故 Sobolev 弱收敛蕴含 \(F_j\rightharpoonup F\) 于 \(L^p(Z)\)。而标准范数恰好满足
\[
 \|F_j\|_{L^p(Z)}=\|u_j\|_{W^{1,p}(\Omega)},\quad
 \|F_j-F\|_{L^p(Z)}=\|u_j-u\|_{W^{1,p}(\Omega)}.
\]
应用前一引理即可。
\end{proof}

这里使用了标准范数的有限分量 \(p\) 次和形式。等价范数保持强、弱拓扑，却不保证“某一范数的数值收敛”对另一个范数仍然成立；不能把上一命题的附加条件换成任意等价范数而不另作证明。

在 \(\Omega=(0,2\pi)\) 上，取
\[
 u_j(x)=\frac{\sin(jx)}j,\quad u_j'(x)=\cos(jx).
\]
函数一致趋于零，由零端点及\cref{b1:thm:interval-zero-boundary}也知它们属于 \(W_0^{1,p}(\Omega)\)，但
\[
 \|u_j'\|_{L^p(0,2\pi)}^p
 =\int_0^{2\pi}|\cos x|^p\,\mathrm dx>0.
\]
同时梯度弱趋于零：对 \(\psi\in C_c^\infty(0,2\pi)\)，分部积分给出
\[
 \left|\int_0^{2\pi}\psi(x)\cos(jx)\,\mathrm dx\right|
 =\frac1j\left|\int_0^{2\pi}\psi'(x)\sin(jx)\,\mathrm dx\right|
 \le\frac{\|\psi'\|_1}{j}\longrightarrow0.
\]
再用这类函数在 \(L^{p'}\) 中的稠密性和余弦函数的统一 \(L^p\) 界，把结论推广到全部测试函数。于是 \(u_j\rightharpoonup0\) 于 \(W^{1,p}\)，却不强收敛。这个例子中的障碍是梯度振荡；函数值的强收敛没有消除导数的范数。

\subsection{变分问题中的应用}
\label{b1:subsec:250304}

下面在一个具体能量上实施\secref{b1:sec:0905}的变分直接法。零阶正项同时控制函数大小，因此不需要借助 Poincaré 不等式，也不必为区域增加边界正则性。

\begin{theorem}\label{b1:thm:sobolev-variational-minimum}
设 \(\Omega\subset\mathbb R^d\) 为有界开集，\(1<p<\infty\)，
\(g\in W^{1,p}(\Omega)\)、\(f\in L^{p'}(\Omega)\) 均为实值。则泛函
\[
 J(u)=\frac1p\int_\Omega\bigl(|\nabla u|^p+|u|^p\bigr)\,\mathrm dx
                -\int_\Omega fu\,\mathrm dx
\]
在仿射集 \(\mathcal A=g+W_0^{1,p}(\Omega)\) 上有唯一极小点 \(u\)。这里
\(|\nabla u|\) 取欧氏长度。极小点也恰是 \(\mathcal A\) 中满足下列积分恒等式的函数：
\begin{equation}\label{b1:eq:sobolev-euler-lagrange}
 \int_\Omega\left(
   |\nabla u|^{p-2}\nabla u\cdot\nabla\varphi
       +|u|^{p-2}u\varphi
 \right)\,\mathrm dx
 =\int_\Omega f\varphi\,\mathrm dx
 \quad\bigl(\varphi\in W_0^{1,p}(\Omega)\bigr).
\end{equation}
其中 \(|z|^{p-2}z\) 在 \(z=0\) 处规定为零。
\end{theorem}
\begin{proof}
记
\[
 A(u)=\int_\Omega\bigl(|\nabla u|^p+|u|^p\bigr)\,\mathrm dx.
\]
有限维范数的等价性给出 \(c_{d,p}\|u\|_{W^{1,p}}^p\le A(u)\)。当 \(p\ne2\) 时，不把这里的欧氏梯度能量与逐分量 \(p\) 次和认作数值相等。由 Hölder 不等式及带参数的 Young 不等式，
\[
 \left|\int_\Omega fu\,\mathrm dx\right|
 \le\|f\|_{p'}\cdot\|u\|_p
 \le\frac1{2p}\|u\|_p^p+C_p\|f\|_{p'}^{p'}.
\]
所以
\[
 J(u)\ge\frac1{2p}A(u)-C_p\|f\|_{p'}^{p'}
 \ge c\|u\|_{W^{1,p}}^p-C_p\|f\|_{p'}^{p'}.
\]
这既给出下界，也保证能量有统一上界的序列在 \(W^{1,p}\) 中有界。又 \(\mathcal A\) 非空且 \(J(g)\) 有限，故其下确界是实数，可以选取极小化序列。

再看弱下半连续性。映射 \(z\mapsto|z|^p\) 在欧氏空间中凸，因此 \(A\) 凸。它也在 Sobolev 范数下连续：梯度的欧氏 \(L^p\) 范数之差不超过梯度差的欧氏 \(L^p\) 范数，后者由标准 Sobolev 范数控制；零阶项相同。取 \(p\) 次幂仍保持连续。由\cref{b1:thm:convex-weak-lsc}，\(A\) 弱下半连续，而 \(u\mapsto\int fu\) 弱连续。因此 \(J\) 弱下半连续。

从极小化序列中抽出 Sobolev 弱收敛子列。由\cref{b1:prop:sobolev-weak-limit-identification}，极限仍在 \(\mathcal A\) 中；由弱下半连续性，它的能量不超过原下确界，因而取得极小值。若两个实函数代表不同的等价类，它们在正测度集合上不同，严格凸的零阶函数 \(s\mapsto|s|^p\) 使中点能量严格小于两个能量的平均。梯度项凸、线性项不影响这个严格不等式，故极小点唯一。

最后推导方程，并核实积分求导。对固定 \(\varphi\in W_0^{1,p}(\Omega)\)，所有实数 \(t\) 都满足 \(u+t\varphi\in\mathcal A\)。欧氏函数 \(z\mapsto|z|^p/p\) 在包括零点在内的每一点可微，导数为 \(|z|^{p-2}z\)。对 \(0<|t|\le1\)，沿线段积分其导数，得到
\[
 \left|\frac{|z+tw|^p-|z|^p}{pt}\right|
 \le C_p\bigl(|z|^{p-1}+|w|^{p-1}\bigr)\cdot|w|.
\]
分别取 \((z,w)=(\nabla u,\nabla\varphi)\) 与 \((u,\varphi)\)。右端由 Hölder 不等式可积，故控制收敛定理允许在两个能量积分内通过差商极限。由
\[
 \left.\frac{\mathrm d}{\mathrm dt}J(u+t\varphi)\right|_{t=0}=0
\]
即得\eqref{b1:eq:sobolev-euler-lagrange}。

反过来，若 \(u\in\mathcal A\) 满足该恒等式，对任意 \(v\in\mathcal A\) 取 \(\varphi=v-u\)。凸函数的一阶支撑不等式给出
\[
 \begin{split}
 J(v)-J(u)\ge
 \int_\Omega\bigl(
  |\nabla u|^{p-2}\nabla u\cdot\nabla(v-u)
       +|u|^{p-2}u(v-u)-f(v-u)
 \bigr)\,\mathrm dx=0.
 \end{split}
\]
因此 \(u\) 是极小点，且由已经证明的严格凸性唯一。
\end{proof}

式\eqref{b1:eq:sobolev-euler-lagrange}称为这一能量的%
\term*[euler-lagrange-fangcheng]{Euler–Lagrange 方程}[Euler–Lagrange equation]的弱形式。它在推导时只要求一阶弱导数，没有预先假定极小点具有二阶经典导数。

这个存在性证明没有使用 Rellich 紧性：自反性给出弱子列，闭仿射约束保留边界条件，凸性给出能量的弱下半连续性。若再加入关于 \(u\) 的非凸低阶项，弱下半连续性可能失效，此时紧嵌入提供的 \(L^q\) 强收敛才可能用来控制该项，还须核对它的连续性和增长条件。不能仅因已经抽取了弱子列，就宣称全部非线性项都能通过极限。

希望把这些步骤放在 Sobolev 空间与偏微分方程的系统课程中学习，可伴读 Brezis 的《Functional Analysis, Sobolev Spaces and Partial Differential Equations》\cite{Brezis2011Functional}。本节的极小值论证也说明，使用紧性之前先辨认方程中究竟需要哪一种收敛，往往能够减少不必要的几何假设。


\begin{chaptersummary}
有限指数 \(L^p\) 的强紧性可以落实为一个具体逼近过程：先用小平移控制网格内的变化，再用尾部控制把无穷多个网格缩减为有限多个。有界弱梯度提供了整族的一致平移模；有界区域及适当延拓提供尾部控制，由此得到最基本的 Sobolev 紧嵌入。更高积分指数和较低 Hölder 指数的紧性，通过连续估计与插值进一步获得。

边界规则性不是所有版本都需要的条件。一般函数可借 Lipschitz 延拓，零边界闭包可以补零，局部问题可以先截断；高阶问题则必须有真正控制相应导数的延拓。临界指数处不再保留插值余量，集中例子说明这会实质性地破坏紧性。

弱收敛保留线性导数关系与闭凸约束，紧性能够加强像空间中的收敛，而原 Sobolev 范数的强收敛还需要额外信息。凸能量的极小值问题可以仅用弱下半连续性；需要保留单位范数等非弱闭约束，或处理合适的非凸低阶项时，强紧性才承担不可替代的任务。下一章进一步研究等价类在多小的例外集外能够确定点值，这与本章的积分范数紧性是另一层问题。
\end{chaptersummary}

\BookExercises

% \chapref{b1:ch:25}习题临时稿；不另设 BookExercises 入口。
% 八题均按本章及前章已建立的工具自拟，来源与实质增量逐题说明，不冒认教材原题号。
% 不使用容量、谱完备性或尚未建立的偏微分方程正则性。

\begin{exercise}\label{b1:ex:25-fk-redundant-boundedness}
设 \(1\le p<\infty\)，\(\mathcal F\subset L^p(\mathbb R^d)\) 满足
\[
 \lim_{R\to\infty}\sup_{f\in\mathcal F}
       \|f\|_{L^p(\{|x|>R\})}=0,
 \quad
 \lim_{\delta\downarrow0}\sup_{f\in\mathcal F}
       \sup_{|h|<\delta}\|\tau_hf-f\|_p=0,
\]
其中 \(\tau_hf(x)=f(x-h)\)。证明 \(\mathcal F\) 的 \(L^p\) 范数自动一致有界。
因此，Fréchet–Kolmogorov 判据虽然可以列成三个条件，却不能据此断言三个条件在逻辑上相互独立。
\end{exercise}
% 来源：自拟推导；用长位移的有限分步把球内质量与球外尾部相比较，不借用未实读论文的证明。
\begin{solution}
取 \(R\ge1\)，使每个 \(f\in\mathcal F\) 的球外范数不超过 \(1\)，再取
\(\delta>0\)，使所有 \(|h|<\delta\) 的平移误差都不超过 \(1\)。固定
\(a=3Re_1\)，选一个正整数 \(N\)，使 \(|a|/N<\delta\)，并记 \(h=a/N\)。
平移等距性与望远镜求和给出
\[
 \tau_af-f=\sum_{j=0}^{N-1}\tau_{jh}(\tau_hf-f),
 \quad
 \|\tau_af-f\|_p\le N.
\]
对 \(x\in B(0,R)\)，有 \(|x-a|>R\)，因此
\[
 \|\tau_af\|_{L^p(B(0,R))}
   =\|f\|_{L^p(B(0,R)-a)}\le1.
\]
于是
\[
 \begin{aligned}
 \|f\|_p
 &\le\|f\|_{L^p(B(0,R))}
              +\|f\|_{L^p(\{|x|>R\})}\\
 &\le\|f-\tau_af\|_p
        +\|\tau_af\|_{L^p(B(0,R))}+1
 \le N+2.
 \end{aligned}
\]
球面为零测集，开闭球的差别不影响这些范数。数 \(N\) 只由两个统一条件选出的
\(R,\delta\) 决定，不依赖 \(f\)，故得到一致界。再结合定理%
\ref{b1:thm:frechet-kolmogorov-lp}，这两个条件本身已经足以推出相对紧性。
\end{solution}

\begin{exercise}\label{b1:ex:25-weighted-confinement}
设 \(1\le p<\infty\)、\(a>0\)、\(0\le M<\infty\)。证明集合
\[
 \mathcal K=\left\{u\in W^{1,p}(\mathbb R^d):
   \|u\|_p^p+\|\nabla u\|_p^p+
        \int_{\mathbb R^d}|x|^a\cdot|u(x)|^p\,\mathrm dx
                               \le M^p\right\}
\]
在 \(L^p(\mathbb R^d)\) 中相对紧。解释最后的加权项如何替代有界区域或共同紧支集的条件。
\end{exercise}
% 来源：自拟；全空间的矩控制给出尾部，弱导数给出一致平移控制，组合为非有界区域上的紧性。
\begin{solution}
集合中的函数首先满足 \(\|u\|_p\le M\)。对 \(R>0\)，由
\(|x|^a\ge R^a\) 在球外成立，得到
\[
 \int_{|x|>R}|u(x)|^p\,\mathrm dx
 \le R^{-a}\int_{\mathbb R^d}|x|^a\cdot|u(x)|^p\,\mathrm dx
 \le M^pR^{-a}.
\]
因此尾部一致趋于零。

还需控制平移。若 \(u\in C_c^\infty(\mathbb R^d)\)，线段积分给出
\[
 u(x-h)-u(x)
       =-\int_0^1\nabla u(x-th)\cdot h\,\mathrm dt,
\]
从而 Minkowski 积分不等式及平移等距性推出
\[
 \|\tau_hu-u\|_p\le|h|\cdot\|\nabla u\|_p.
\]
由\cref{b1:thm:whole-space-sobolev-density}，对一般
\(u\in W^{1,p}(\mathbb R^d)\) 作光滑紧支集逼近；等式两侧所需的
\(L^p\) 范数都可通过极限，上式仍成立。因此对 \(\mathcal K\) 有
\[
 \sup_{u\in\mathcal K}\|\tau_hu-u\|_p\le M|h|\longrightarrow0.
\]
定理\ref{b1:thm:frechet-kolmogorov-lp}的三个条件均已满足，所求相对紧性随之成立。

加权项并不要求任何一个函数具有紧支集；它要求所有函数在远处放置的质量都受到同一个界限制。导数界约束局部振荡，而这个额外条件阻止质量整体移向无穷远，两者承担不同的任务。
\end{solution}

\begin{exercise}\label{b1:ex:25-critical-mass-concentration}
设 \(d\ge2\)、\(1\le p<d\)、\(p^*=dp/(d-p)\)，取非零
\(\varphi\in C_c^\infty(B(0,1))\)。对 \(0<\varepsilon<1\)，令
\[
 u_\varepsilon(x)=\varepsilon^{1-d/p}\varphi(x/\varepsilon),
 \quad x\in\mathbb R^d.
\]
\begin{enumerate}
 \item\label{b1:item:25-critical-no-strong-subsequence}
 证明 \(u_\varepsilon\) 在 \(W_0^{1,p}(B(0,1))\) 中有界，在每个
 \(1\le q<p^*\) 的 \(L^q\) 中趋于零；但对于任意 \(\varepsilon_n\downarrow0\)，
 \((u_{\varepsilon_n})\) 都没有 \(L^{p^*}\) 强收敛子列。
 \item\label{b1:item:25-critical-defect-measures}
 把
 \[
  \mu_\varepsilon=|u_\varepsilon|^{p^*}\,\mathrm dx,\quad
  \nu_\varepsilon=|\nabla u_\varepsilon|^p\,\mathrm dx
 \]
 看成全空间上的有限正 Borel 测度。证明
 \[
  \mu_\varepsilon\Rightarrow
      \|\varphi\|_{p^*}^{p^*}\delta_0,
  \quad
  \nu_\varepsilon\Rightarrow
      \|\nabla\varphi\|_p^p\delta_0.
 \]
\end{enumerate}
\end{exercise}
% 来源：自拟；临界重缩放后进一步计算有限测度的弱极限，不止于前章排除过强范数估计。
\begin{solution}
每个函数光滑且具有包含在单位球内部的紧支集。变量代换给出
\[
 \|\nabla u_\varepsilon\|_p=\|\nabla\varphi\|_p,
 \quad
 \|u_\varepsilon\|_q
  =\varepsilon^{1-d/p+d/q}\|\varphi\|_q
                      \quad(1\le q<\infty).
\]
特别地，\(\|u_\varepsilon\|_p=\varepsilon\|\varphi\|_p\)，所以
\(W^{1,p}\) 范数一致有界。当 \(q<p^*\) 时，指数
\(1-d/p+d/q\) 为正，得到强 \(L^q\) 收敛；在 \(q=p^*\) 时，这个指数等于零，因而
\[
 \|u_\varepsilon\|_{p^*}=\|\varphi\|_{p^*}>0.
\]
对每个 \(x\ne0\)，当 \(\varepsilon\) 足够小时，\(x\) 已在
\(\varepsilon\operatorname{supp}\varphi\) 之外，所以 \(u_\varepsilon(x)=0\)。
假如某个子列在 \(L^{p^*}\) 中强收敛，则它还能抽出几乎处处收敛的子列；上述逐点结论使其极限只能为零。但强收敛应使范数趋于零，与固定的正范数矛盾。这排除了全部强收敛子列。

对第二问，任取 \(\psi\in C_b(\mathbb R^d)\)。临界指数恰好消去变量代换产生的尺度因子，故
\[
 \int\psi\,\mathrm d\mu_\varepsilon
    =\int_{\mathbb R^d}\psi(\varepsilon y)|\varphi(y)|^{p^*}\,\mathrm dy
    \longrightarrow\psi(0)\|\varphi\|_{p^*}^{p^*}.
\]
控制收敛适用，因为 \(\psi\) 有界，而 \(|\varphi|^{p^*}\) 可积。同样，
\[
 \int\psi\,\mathrm d\nu_\varepsilon
    =\int_{\mathbb R^d}\psi(\varepsilon y)|\nabla\varphi(y)|^p\,\mathrm dy
    \longrightarrow\psi(0)\|\nabla\varphi\|_p^p.
\]
这就是有限测度弱收敛的定义。函数在次临界范数中消失，并不意味着这些临界总量消失；它们集中到了一个点上。
\end{solution}

\begin{exercise}\label{b1:ex:25-holder-endpoint}
设 \(d<p<\infty\)、\(\alpha=1-d/p\)，取
\(\varphi\in C_c^\infty(B(0,1/4))\) 且 \(\varphi(0)=1\)。对
\(0<\varepsilon<1\)，在单位球 \(B\) 上令
\[
 u_\varepsilon(x)=\varepsilon^\alpha\varphi(x/\varepsilon).
\]
证明这是一族有界的 \(W_0^{1,p}(B)\) 函数，并且对每个
\(0<\beta<\alpha\)，有 \(u_\varepsilon\to0\) 于
\(C^{0,\beta}(\overline B)\)。另一方面，证明对任何
\(\varepsilon_n\downarrow0\)，它都没有 \(C^{0,\alpha}(\overline B)\) 强收敛子列。
由此区分 Morrey 的连续嵌入指数与 Rellich 的紧嵌入指数。
\end{exercise}
% 来源：自拟；在可用的临界 Hölder 估计内证明失去相对紧性，不再讨论更高指数是否有统一界。
\begin{solution}
与上一题相同的变量代换给出
\[
 \|u_\varepsilon\|_p=\varepsilon\|\varphi\|_p,
 \quad
 \|\nabla u_\varepsilon\|_p=\|\nabla\varphi\|_p,
 \quad
 \|u_\varepsilon\|_\infty=\varepsilon^\alpha\|\varphi\|_\infty.
\]
函数均光滑且内部紧支集，所以属于 \(W_0^{1,p}(B)\)，并在其中有界；最后一式还给出一致趋零。

把 \(\varphi\) 视为全空间光滑紧支集函数。它的全空间
\(\beta\)-Hölder 半范数有限：距离不超过 \(1\) 时用梯度上界，距离超过 \(1\) 时用两倍上确界。对任意 \(x,y\in B\) 作缩放比较，得到
\[
 [u_\varepsilon]_{0,\beta;\overline B}
     \le\varepsilon^{\alpha-\beta}
                   [\varphi]_{0,\beta;\mathbb R^d}.
\]
若 \(\beta<\alpha\)，这个半范数也趋于零，因而得到所求的强
\(C^{0,\beta}\) 收敛。

但取 \(a=e_1/2\)，有 \(\varphi(a)=0\)。在 \(0,\varepsilon a\) 两点比较，
\[
 [u_\varepsilon]_{0,\alpha;\overline B}
 \ge\frac{|u_\varepsilon(0)-u_\varepsilon(\varepsilon a)|}
             {|\varepsilon a|^\alpha}
 =2^\alpha.
\]
若某个子列在 \(C^{0,\alpha}\) 中强收敛，它也应一致收敛；唯一可能的极限是零。
然而强 \(C^{0,\alpha}\) 收敛到零又要求上述半范数趋于零，矛盾。
所以临界指数处有连续嵌入，并不足以推出同一目标空间中的相对紧性。
\end{solution}

\begin{exercise}\label{b1:ex:25-infinitely-many-components}
设 \(d\ge1\)、\(1\le p<\infty\)，令
\[
 x_n=2^{-n}e_1,\quad r_n=2^{-n-3},\quad
 B_n=B(x_n,r_n),\quad
 \Omega=\bigcup_{n=1}^\infty B_n.
\]
证明 \(\Omega\) 是有界开集，但嵌入
\(W^{1,p}(\Omega)\to L^p(\Omega)\) 不紧。具体地，考虑
\[
 u_n=|B_n|^{-1/p}\mathbf1_{B_n}
                  \quad\text{作为 }\Omega\text{ 上的函数},
\]
计算它们的弱导数、Sobolev 范数和两两 \(L^p\) 距离。
解释这为什么不反驳任意有界开集上的 \(W_0^{1,p}\) 紧嵌入。
\end{exercise}
% 来源：自拟；无限分量上的归一化常数，区别无边界条件空间与零边界闭包。
\begin{solution}
相邻两个中心的距离为 \(2^{-n-1}\)，而半径之和为
\(3\cdot2^{-n-4}\)，前者严格较大；不相邻的球距离更大。因此各球闭包两两不交。
所有球都包含在单位球内，故 \(\Omega\) 是有界开集。中心列趋于原点，而原点不属于 \(\Omega\)。

函数 \(u_n\) 在每个分量内都是常数，因此其弱梯度在 \(\Omega\) 内为零。也可以直接检验：若 \(\phi\in C_c^\infty(\Omega)\)，则
\(\phi|_{B_n}\) 的支集与 \(\partial B_n\) 保持正距离，故对每个坐标 \(i\)，
\[
 \int_\Omega u_n\partial_i\phi
      =|B_n|^{-1/p}\int_{B_n}\partial_i\phi=0.
\]
这里是在 \(\Omega\) 内求弱导数，没有把 \(u_n\) 补零后跨越分量边界求导。
由归一化可得
\[
 \|u_n\|_{L^p(\Omega)}=1,\quad
 \nabla u_n=0,\quad
 \|u_n\|_{W^{1,p}(\Omega)}=1.
\]
当 \(n\ne m\) 时，两个函数支集不交，因此
\[
 \|u_n-u_m\|_{L^p(\Omega)}^p
          =\|u_n\|_p^p+\|u_m\|_p^p=2.
\]
这是一列有界 Sobolev 函数，却没有 \(L^p\) Cauchy 子列，故嵌入不紧。

最后，任何一个 \(u_n\) 都不属于 \(W_0^{1,p}(\Omega)\)。否则定理%
\ref{b1:thm:zero-boundary-poincare}会使
\(\|u_n\|_p\le C\|\nabla u_n\|_p=0\)，与其范数等于 \(1\) 矛盾。
因此这个例子检验的是不加边界条件且不加边界规则性时的失败，不涉及零边界版本的反例。
\end{solution}

\begin{exercise}\label{b1:ex:25-nonlinear-reaction-limit}
设 \(\Omega\subset\mathbb R^d\) 为有界 Lipschitz 开集，
\(1\le p<d\)、\(1<q<p^*=dp/(d-p)\)，并令 \(q'=q/(q-1)\)。
设实值函数 \(u_n\rightharpoonup u\) 于 \(W^{1,p}(\Omega)\)。
\begin{enumerate}
 \item\label{b1:item:25-reaction-strong-limit}
 证明
 \[
  |u_n|^{q-2}u_n\longrightarrow |u|^{q-2}u
                       \quad\text{于 }L^{q'}(\Omega),
 \]
 其中在函数值为零处规定相应表达式为零。分别给出 \(q\ge2\) 和
 \(1<q<2\) 所需的点态估计。
 \item\label{b1:item:25-reaction-strong-weak-pairing}
 若另有 \(v_n\rightharpoonup v\) 于 \(L^q(\Omega)\)，证明
 \[
  \int_\Omega |u_n|^{q-2}u_n v_n
          \longrightarrow\int_\Omega |u|^{q-2}u v.
 \]
\end{enumerate}
\end{exercise}
% 来源：自拟；次临界紧嵌入与幂次映射的定量连续性，再组合强弱配对，不把非线性操作直接套在弱极限上。
\begin{solution}
由\cref{b1:thm:rellich-subcritical}，嵌入
\(W^{1,p}(\Omega)\to L^q(\Omega)\) 是紧线性算子。命题%
\ref{b1:prop:compact-weak-to-strong}因而给出整列
\(u_n\to u\) 于 \(L^q(\Omega)\)，不需要再为不同测试函数分别抽子列。
记 \(F(t)=|t|^{q-2}t\)。首先
\[
 \|F(w)\|_{q'}=\|w\|_q^{q-1},
\]
所以题中的非线性项都属于所需对偶指数空间。

当 \(q>2\) 时，由 \(F'(t)=(q-1)|t|^{q-2}\) 及一维积分公式，
\[
 |F(a)-F(b)|
       \le(q-1)(|a|+|b|)^{q-2}\cdot|a-b|.
\]
以指数 \(q/(q-2)\) 和 \(q\) 对乘积使用 Hölder 不等式，得到
\[
 \|F(u_n)-F(u)\|_{q'}
 \le(q-1)(\|u_n\|_q+\|u\|_q)^{q-2}
                              \cdot\|u_n-u\|_q
 \longrightarrow0.
\]
在 \(q=2\) 时，\(F(t)=t\)，结论直接就是 \(L^2\) 强收敛。

当 \(1<q<2\) 时，记 \(\theta=q-1\in(0,1)\)。同号的 \(a,b\) 可用
\(t^\theta\) 的次可加性得到
\(|F(a)-F(b)|\le|a-b|^\theta\)；异号时，由凹性有
\[
 |F(a)-F(b)|
   =|a|^\theta+|b|^\theta
   \le2^{1-\theta}(|a|+|b|)^\theta
   =2^{1-\theta}|a-b|^\theta.
\]
因而对所有实数都有
\[
 |F(a)-F(b)|\le2^{2-q}|a-b|^{q-1},
 \quad
 \|F(u_n)-F(u)\|_{q'}
                  \le2^{2-q}\|u_n-u\|_q^{q-1}\longrightarrow0.
\]
这完成第一问。

对于第二问，弱收敛列 \((v_n)\) 的 \(L^q\) 范数一致有界。分解为
\[
 \int_\Omega F(u_n)v_n-\int_\Omega F(u)v
   =\int_\Omega\bigl(F(u_n)-F(u)\bigr)v_n
                       +\int_\Omega F(u)(v_n-v).
\]
第一项的绝对值不超过
\(\|F(u_n)-F(u)\|_{q'}\cdot\|v_n\|_q\)，所以趋于零；第二项则因
\(F(u)\in L^{q'}\) 是一个固定的弱收敛测试函数而趋于零。
这里先由紧性恢复了 \(u_n\) 的强收敛，才有资格把非线性项传到极限。
\end{solution}

\begin{exercise}\label{b1:ex:25-first-dirichlet-eigenvalue}
设 \(\Omega\subset\mathbb R^d\) 为非空有界开集，本题取实值函数，并定义
\[
 \lambda_1=\inf\left\{
       \int_\Omega|\nabla v|^2:
       v\in H_0^1(\Omega),\ \|v\|_{L^2(\Omega)}=1
                    \right\}.
\]
\begin{enumerate}
 \item\label{b1:item:25-first-eigenvalue-attainment}
 证明 \(0<\lambda_1<\infty\)，且存在 \(u\in H_0^1(\Omega)\) 取得这个下确界。
 \item\label{b1:item:25-first-eigenvalue-equation}
 不借助抽象乘子定理，沿保持约束的曲线证明
 \[
  \int_\Omega\nabla u\cdot\nabla\phi
       =\lambda_1\int_\Omega u\phi
              \quad\bigl(\phi\in H_0^1(\Omega)\bigr).
 \]
 若非零 \(w\in H_0^1(\Omega)\) 也满足这个恒等式，但参数换成实数
 \(\lambda\)，证明 \(\lambda\ge\lambda_1\)。
\end{enumerate}
本题的方程按弱导数解释，不要求证明特征函数的经典光滑性或全部特征函数的完备性。
\end{exercise}
% 来源：自拟；L2 单位球这一非凸约束下的直接法与归一化变分，不重复正文带正零阶项的无约束凸能量问题。
\begin{solution}
由\cref{b1:thm:zero-boundary-poincare}，存在 \(C_P<\infty\)，使
\(\|v\|_2\le C_P\|\nabla v\|_2\) 对所有 \(v\in H_0^1(\Omega)\) 成立。
于是约束集中的每个函数都有梯度能量至少 \(C_P^{-2}\)，故
\(\lambda_1>0\)。又因为非空开集包含一个球，可在其中取非零光滑紧支集函数并归一化，得到约束集非空且 \(\lambda_1<\infty\)。

取极小化序列 \(u_n\)，满足 \(\|u_n\|_2=1\) 及
\(\|\nabla u_n\|_2^2\to\lambda_1\)。这列函数在 Hilbert 空间
\(H_0^1(\Omega)\) 中有界。由\cref{b1:thm:reflexive-weak-subsequence}，可以抽出一个子列在该空间中弱收敛到某个 \(u\)。
定理\ref{b1:thm:rellich-subcritical}与\ref{b1:thm:rellich-holder}的零边界版本给出
\(H_0^1(\Omega)\to L^2(\Omega)\) 的紧性：在 \(d>2\) 时
\(2<2d/(d-2)\)，在 \(d=2\) 时使用临界的有限指数结论；在 \(d=1\) 时，紧 Hölder 嵌入先给出一致收敛，再由区域有界性推出 \(L^2\) 收敛。
结合命题\ref{b1:prop:compact-weak-to-strong}，该子列在 \(L^2\) 中强收敛到
\(u\)，所以 \(\|u\|_2=1\)。

梯度范数的弱下半连续性给出
\[
 \|\nabla u\|_2^2
      \le\liminf_n\|\nabla u_n\|_2^2=\lambda_1.
\]
由于 \(u\) 已满足约束，下确界的定义又给出反向不等式，因此
\(\|\nabla u\|_2^2=\lambda_1\)。强 \(L^2\) 收敛在这里保住了单位范数约束；仅有弱收敛并不能保证这一点。

对任意 \(\phi\in H_0^1(\Omega)\)，令
\[
 v_t=\frac{u+t\phi}{\|u+t\phi\|_2}
\]
并取充分小的实数 \(t\)，使分母非零。这个曲线始终属于约束集。
记 \(A=\int_\Omega\nabla u\cdot\nabla\phi\)、\(B=\int_\Omega u\phi\)，则
\[
 \int_\Omega|\nabla v_t|^2
 =\frac{\lambda_1+2tA+t^2\|\nabla\phi\|_2^2}
        {1+2tB+t^2\|\phi\|_2^2}.
\]
左侧在 \(t=0\) 取得极小值，故其导数为零，即
\(2A-2\lambda_1B=0\)。这正是所求恒等式。特别地，对光滑紧支集测试函数成立，所以它表示分布意义的
\(-\Delta u=\lambda_1u\)，零边界条件则由 \(u\in H_0^1(\Omega)\) 表达。

最后，若非零 \(w\) 满足参数为 \(\lambda\) 的同类弱方程，取测试函数
\(\phi=w\)，便有
\[
 \lambda=\frac{\|\nabla w\|_2^2}{\|w\|_2^2}\ge\lambda_1.
\]
因此取得的参数确实是所有这种弱 Dirichlet 特征值中的最小值；这个结论不需要预先建立谱分解。
\end{solution}

\begin{exercise}\label{b1:ex:25-recover-weak-limit}
设 \(X\) 为自反 Banach 空间，\(Y\) 为 Banach 空间，
\(T:X\to Y\) 为有界线性单射。设 \((x_n)\) 在 \(X\) 中有界，且
\(Tx_n\to y\) 于 \(Y\) 的范数。
\begin{enumerate}
 \item\label{b1:item:25-reflexive-recovery}
 证明存在唯一 \(x\in X\)，使 \(Tx=y\)，并且整列
 \(x_n\rightharpoonup x\) 于 \(X\)。注意本方向没有假设 \(T\) 为紧算子。
 \item\label{b1:item:25-nonreflexive-recovery-failure}
 证明去掉 \(X\) 的自反性后，结论即使对紧线性单射也会失败。具体检验
 \[
  T:\ell^1\longrightarrow\ell^2,\quad
  T(a_1,a_2,\ldots)=(a_1,a_2/2,a_3/3,\ldots)
 \]
 及 \(\ell^1\) 的标准单位向量列。
\end{enumerate}
\end{exercise}
% 来源：自拟；由较弱空间中的强极限识别原空间弱极限，配合紧对角算子展示自反性不可由紧性替代。
\begin{solution}
由\cref{b1:thm:reflexive-weak-subsequence}，有界列可以抽出弱收敛子列
\(x_{n_j}\rightharpoonup x\)。有界线性算子保持弱收敛，因为对每个
\(f\in Y^*\)，\(f\circ T\in X^*\)，所以 \(Tx_{n_j}\rightharpoonup Tx\)。
另一方面，它已在范数中趋于 \(y\)，故也弱收敛到 \(y\)。弱极限唯一，得到
\(Tx=y\)。单射性使这样的 \(x\) 至多有一个。

还须从一个子列推进到整列。假设整列不弱收敛到 \(x\)，则存在
\(g\in X^*\)、\(\varepsilon>0\) 及一个子列，满足
\[
 |g(x_{n_j})-g(x)|\ge\varepsilon\quad\text{对每个 }j.
\]
这个子列仍然有界，再用自反性抽出弱收敛子列，记其弱极限为 \(z\)。
按前面的论证，必有 \(Tz=y\)，从而单射性给出 \(z=x\)。这与所列标量距离始终不小于 \(\varepsilon\) 矛盾，因此整列确实弱收敛。

对第二问，所给对角算子线性且单射，并满足
\(\|Ta\|_2\le\|a\|_1\)。记 \(T_Na\) 为只保留前 \(N\) 个坐标的部分，则
\[
 \|(T-T_N)a\|_2
       \le\frac1{N+1}\|a\|_2
       \le\frac1{N+1}\|a\|_1.
\]
对 \(\ell^1\) 单位球，前 \(N\) 个坐标的像是有限维有界集，具有有限的任意精度网；尾项又由上式一致控制。因此 \(T\) 的单位球像全有界，在完备空间
\(\ell^2\) 中相对紧，说明 \(T\) 是紧算子。

可是对标准单位向量 \(e_n\)，有
\[
 \|e_n\|_1=1,\quad \|Te_n\|_2=1/n\longrightarrow0.
\]
若原向量列弱收敛，其极限由 \(T\) 的连续性和单射性只能为零。但连续线性泛函
\[
 g(a)=\sum_{j=1}^\infty a_j,\quad |g(a)|\le\|a\|_1,
\]
始终满足 \(g(e_n)=1\)，故这个列并不弱收敛到零。紧性能够控制像空间中的收敛，不能替代原空间中抽取弱子列所需的条件。
\end{solution}
